\chapter{बल और जड़त्व का सिद्धांत}\label{ch:g10:inertia}

पैडल चलाना बंद कीजिए और साइकिल लुढ़कते-लुढ़कते रुक जाती है; खींचना बंद
कीजिए और स्लेज वहीं खड़ी रह जाती है। अनुभव कान में कहता है कि गति को
लगातार मेहनत की ख़ुराक चाहिए --- और दो हज़ार साल तक भौतिकी भी यही मानती
रही। यह अध्याय बल गढ़ता है, यानी \omterm{def:g10:inertia:force}{न्यूटन} में नापा जाने वाला
एक तीर, और यांत्रिकी का पहला बड़ा नियम बताता है: सचमुच अकेला छोड़ दिया
जाए तो चलता पिंड अपनी गति सदा बनाए रखता है, सीधी और स्थिर।

\section{किसी क्रिया को बल के रूप में गढ़ना}

\begin{definition}[बल]\label{def:g10:inertia:force}
\emph{बल}\index{बल} किसी पिंड पर लगाई गई यांत्रिक क्रिया (धक्का, खिंचाव,
आकर्षण) का प्रतिरूप है। उसकी पहचान इनसे होती है:
\begin{itemize}
  \item उसका \emph{क्रिया बिंदु}\index{क्रिया बिंदु};
  \item उसकी दिशा --- वह रेखा जिसके अनुदिश वह लगता है, और उस रेखा पर
        उसका रुख़;
  \item उसका परिमाण, जिसे \emph{न्यूटन}\index{न्यूटन} (संकेत \unit{N}) में
        नापा जाता है।
\end{itemize}
बल को उसके क्रिया बिंदु से निकलते एक तीर $\vect F$ (यानी एक सदिश, जैसा गणित
के खंड में है) के रूप में खींचा जाता है, जिसकी लंबाई किसी बताए गए पैमाने
पर परिमाण के समानुपाती होती है।
\end{definition}

\begin{example}[बस्ते का भार]\label{ex:g10:inertia:schoolbag}
$m = \qty{4.0}{kg}$ द्रव्यमान के बस्ते को पृथ्वी उसके भार (\cref{ch:g10:universal-gravitation}) से नीचे की ओर
खींचती है: $P = mg = 4.0 \times 9.81 \approx \qty{39}{N}$ --- और एक \omterm{def:g10:inertia:force}{न्यूटन} लगभग एक सेब के भार
जितना है। पैमाने $\qty{1}{cm} \leftrightarrow
\qty{10}{N}$ पर: बस्ते के गुरुत्व केंद्र पर लगा, नीचे की ओर
इशारा करता \qty{3.9}{cm} लंबा एक \omterm{def:g10:universal-gravitation:weight}{ऊर्ध्वाधर} तीर।
\end{example}

\section{बलों की सूची}

\begin{definition}[संस्पर्श और दूरस्थ बल]\label{def:g10:inertia:contact}
\emph{संस्पर्श बल}\index{संस्पर्श बल} वहाँ लगता है जहाँ दो पिंड एक-दूसरे
को छूते हैं। \emph{दूरस्थ बल}\index{दूरस्थ बल} को स्पर्श चाहिए ही नहीं:
इस अध्याय में ऐसा इकलौता बल भार है --- यानी \cref{ch:g10:universal-gravitation} वाला
\omterm{def:g10:universal-gravitation:force}{गुरुत्वाकर्षण} खिंचाव।
\end{definition}

\begin{definition}[आम संस्पर्श बल]\label{def:g10:inertia:inventory}
चार \omterm{def:g10:inertia:contact}{संस्पर्श बल} रोज़मर्रा की अधिकांश स्थितियाँ सँभाल लेते हैं:
\begin{itemize}
  \item किसी आधार की \emph{अभिलंब प्रतिक्रिया}\index{अभिलंब प्रतिक्रिया}
        $\vect R$, जो उसकी सतह के लंबवत होती है और उससे दूर धकेलती है;
  \item तार, रस्सी या ज़ंजीर का \emph{तनाव}\index{तनाव} $\vect T$, जो उसके
        अनुदिश होता है और उसकी ओर खींचता है;
  \item किसी सतह, वायु या पानी का \emph{घर्षण}\index{घर्षण} $\vect f$, जो
        सरकने या गति का विरोध करता है;
  \item किसी इंजन, जेट या हाथ का \emph{प्रणोद}\index{प्रणोद}, जो मेहनत
        की दिशा में धकेलता है।
\end{itemize}
\end{definition}

\begin{omfigure}
\begin{tikzpicture}[scale=0.9, font=\small]
  \draw[fill=gray!15] (-2.4,0) rectangle (2.4,-0.25);
  \draw (-2.0,-0.25) -- (-2.0,-1.5);
  \draw (2.0,-0.25) -- (2.0,-1.5);
  \draw[fill=omExo!25] (-0.9,0) rectangle (0.9,0.5);
  \coordinate (G) at (0,0.25);
  \fill (G) circle (1.5pt);
  \draw[-Stealth, omThm, very thick] (G) -- (0,-1.3)
    node[right] {$\vect P$};
  \draw[-Stealth, omDef, very thick] (G) -- (0,1.8)
    node[right] {$\vect R$};
\end{tikzpicture}

{\small मेज़ पर रखी पुस्तक: भार $\vect P$ और \omterm{def:g10:inertia:inventory}{अभिलंब
प्रतिक्रिया} $\vect R$ --- वही रेखा, बराबर लंबाइयाँ, उलटे रुख़।}
\end{omfigure}

\begin{example}[खींची जाती स्लेज की सूची]\label{ex:g10:inertia:sled}
बर्फ़ पर खींची जाती स्लेज: भार $\vect P$ (दूरस्थ), बर्फ़ की
\omterm{def:g10:inertia:inventory}{अभिलंब प्रतिक्रिया} $\vect R$, रस्सी का \omterm{def:g10:inertia:inventory}{तनाव} $\vect T$, और
सरकने का विरोध करता \omterm{def:g10:inertia:inventory}{घर्षण} $\vect f$। स्लेज को और कुछ छूता नहीं:
सूची पूरी हो गई --- \emph{पहले हमेशा भार, फिर हर छूती हुई वस्तु
के लिए एक \omterm{def:g10:inertia:contact}{संस्पर्श बल}}।
\end{example}

\section{जड़त्व का सिद्धांत}

\begin{definition}[सरल रेखीय एकसमान गति]\label{def:g10:inertia:uniform}
कोई पिंड किसी तंत्र में \emph{सरल रेखीय एकसमान गति}\index{सरल रेखीय
एकसमान गति} में है जब उसका \omterm{def:g10:relative-motion:trajectory}{गतिपथ} एक सरल रेखा हो और वह उस पर अचर
चाल से चले: यानी उसका वेग सदिश $\vect v$ हर क्षण वही रहे। \omterm{def:g10:relative-motion:frame}{विराम} इसकी विशेष
अवस्था $\vect v = \vect 0$ है।
\end{definition}

\begin{definition}[संतुलनकारी बल]\label{def:g10:inertia:compensate}
किसी पिंड पर लगते बल \emph{संतुलनकारी}\index{संतुलनकारी बल} तब
होते हैं जब उनका मिला-जुला असर शून्य हो: नोक-से-पूँछ जोड़ने पर उनके तीर
वहीं लौट आते हैं जहाँ से चले थे। दो बल ठीक तभी संतुलनकारी होते
हैं जब उनकी क्रिया-रेखा एक ही हो, परिमाण बराबर हों और रुख़ उलटे।
\end{definition}

\begin{theorem}[जड़त्व का सिद्धांत]\label{thm:g10:inertia:principle}
\index{जड़त्व का सिद्धांत}\index{जड़त्व}
\omterm{def:g10:relative-motion:frames}{भू-निर्देश तंत्र} में जिस पिंड पर कोई बल न लगे, या
\omterm{def:g10:inertia:compensate}{संतुलनकारी बल} लगें, वह \omterm{def:g10:relative-motion:frame}{विराम} में बना रहता है या अपनी
\omterm{def:g10:inertia:uniform}{सरल रेखीय एकसमान गति} बनाए रखता है; और इसका उलटा भी सही है:
\omterm{def:g10:relative-motion:frame}{विराम} में या \omterm{def:g10:inertia:uniform}{सरल रेखीय एकसमान गति} में रहते पिंड पर
\omterm{def:g10:inertia:compensate}{संतुलनकारी बल} ही लगते हैं (या कोई बल लगता ही नहीं)।
\end{theorem}
\admitted

\begin{remark}
कोई प्रयोग किसी पिंड को पूरी तरह अलग-थलग नहीं कर सकता, इसलिए यह सिद्धांत
एक अभिगृहीत है --- गैलीलियो को इसकी झलक मिली, और न्यूटन ने इसे कहा। आगे
का एक अध्याय इसे उस नियम तक चढ़ा देता है जो बलों को वेग के
\emph{परिवर्तनों} से जोड़ता है; और स्नातक वर्ष 1 का खंड पूछता है कि यह
असल में किन तंत्रों में लागू होता है।
\end{remark}

\begin{proposition}[बलों से गति पढ़ना, और गति से बल]\label{prop:g10:inertia:converse}
\omterm{def:g10:relative-motion:frames}{भू-निर्देश तंत्र} में:
\begin{enumerate}
  \item यदि किसी पिंड का वेग बदलता है --- परिमाण में \emph{या} दिशा में
        --- तो उस पर लगते बल \omterm{def:g10:inertia:compensate}{संतुलनकारी} \emph{नहीं} हैं;
  \item यदि कोई पिंड \omterm{def:g10:relative-motion:frame}{विराम} में है या \omterm{def:g10:inertia:uniform}{सरल रेखीय एकसमान गति} में
        है, तो उसकी सूची के बल आपस में कट जाने ही चाहिए --- और
        इससे अज्ञात राशियाँ निकल आती हैं।
\end{enumerate}
\end{proposition}

\begin{proof}
पहला कथन \cref{thm:g10:inertia:principle} का प्रतिधनात्मक है, और दूसरा उसी के विलोम भाग को पूरी की
गई सूची पर लगाने से मिलता है।
\end{proof}

\begin{omfigure}
\begin{tikzpicture}[scale=0.85, font=\small]
  \draw[gray!60] (-0.5,0) -- (10.3,0) node[below left, gray] {बर्फ़};
  \foreach \x in {0.5,2.5,4.5,6.5,8.5} {
    \fill (\x,0.13) ellipse (0.3 and 0.13);
    \draw[-Stealth, omDef, thick] (\x,0.5) -- (\x+1.3,0.5);
  }
  \node[omDef] at (1.15,0.9) {$\vect v$};
  \node[omDef] at (5.4,1.5) {बराबर समय में बराबर क़दम};
\end{tikzpicture}

{\small छोड़े जाने के बाद हॉकी की चकती: बराबर समय-अंतरालों पर उसकी
स्थितियाँ, और हर स्थिति पर एक जैसा वेग तीर --- \omterm{def:g10:inertia:uniform}{सरल रेखीय एकसमान
गति}।}
\end{omfigure}

\begin{example}[हॉकी की चकती]\label{ex:g10:inertia:puck}
छड़ी से छोड़े जाने के बाद चकती चिकनी बर्फ़ पर \qty{30}{m} की दूरी \qty{2.5}{s} में,
सीधे और अचर $v = 30/2.5 = \qty{12}{m/s}$ से तय करती है। सूची: भार और \omterm{def:g10:inertia:inventory}{अभिलंब
प्रतिक्रिया} (\omterm{def:g10:inertia:inventory}{घर्षण} नगण्य)। ये \omterm{def:g10:inertia:compensate}{संतुलनकारी} हैं, और
सिद्धांत ठीक वही गति बताता है जो देखी जाती है: चकती को आगे कुछ भी धकेल
नहीं रहा --- और धकेलने की ज़रूरत भी नहीं।
\end{example}

\begin{remark}[अरस्तू के विरुद्ध गैलीलियो]
अरस्तू सिखाते थे कि हर गति को कोई चलाने वाला चाहिए, और रोज़मर्रा की
ज़िंदगी --- जहाँ हर संपर्क में \omterm{def:g10:inertia:inventory}{घर्षण} है --- मानो सहमत भी दिखती
है। गैलीलियो ने \omterm{def:g10:inertia:inventory}{घर्षण} को आदर्श रूप में हटा दिया: सतह जितनी अधिक
चिकनी हो, गति \emph{बनाए रखने} के लिए उतना ही कम धक्का चाहिए। स्वाभाविक
अवस्था \omterm{def:g10:relative-motion:frame}{विराम} नहीं, बल्कि गति का बना रहना है।
\end{remark}

\begin{remark}[झटका खाता यात्री]
जब बस अचानक ब्रेक लगाती है तो खड़े यात्री आगे की ओर लुढ़क जाते हैं। कोई
बल उन्हें फेंकता नहीं: जड़त्व के कारण उनके शरीर अपना वेग बनाए
रखते हैं जबकि बस अपना वेग खो देती है। सीट-बेल्ट वही पीछे की ओर लगता
बल जुटाती है जिसकी सूची में कमी थी।
\end{remark}

\begin{omfigure}
\begin{tikzpicture}[scale=0.62, font=\small]
  \draw[gray!60] (-0.3,0) -- (7.3,0);
  \foreach \x in {0.7,2.5,4.3,6.1}
    \fill[omDef] (\x,0.22) ellipse (0.34 and 0.22);
  \draw[-Stealth, omThm, thick] (0.7,0.22) -- (0.7,-1.0)
    node[right] {$\vect P$};
  \draw[-Stealth, omDef, thick] (0.7,0.44) -- (0.7,1.65)
    node[right] {$\vect R$};
  \node at (3.6,-1.7) {बल संतुलित: बराबर क़दम};
\end{tikzpicture}
\qquad
\begin{tikzpicture}[scale=0.62, font=\small]
  \draw[gray!60] (-0.3,0) -- (7.3,0);
  \foreach \x in {0.7,2.4,3.8,4.9,5.7,6.2}
    \fill[omProp] (\x,0.22) ellipse (0.34 and 0.22);
  \draw[-Stealth, omThm, thick] (0.7,0.22) -- (0.7,-1.0)
    node[right] {$\vect P$};
  \draw[-Stealth, omDef, thick] (0.7,0.44) -- (0.7,1.65)
    node[right] {$\vect R$};
  \draw[-Stealth, omProp, very thick] (0.7,0.22) -- (-0.3,0.22)
    node[above] {$\vect f$};
  \node at (3.6,-1.7) {बिना कटा घर्षण: पत्थर धीमा पड़ता है};
\end{tikzpicture}

{\small बराबर समय-अंतरालों पर कर्लिंग का पत्थर। बाएँ: $\vect P$ और $\vect R$
\omterm{def:g10:inertia:compensate}{संतुलनकारी} हैं --- क़दम बराबर। दाएँ: एक बिना कटा \omterm{def:g10:inertia:inventory}{घर्षण}
$\vect f$ बचा रह जाता है --- क़दम सिकुड़ते जाते हैं, वेग बदलता जाता है।}
\end{omfigure}

\section{संतुलित बल: आरेख से स्थैतिकी}

\begin{definition}[बल आरेख]\label{def:g10:inertia:diagram}
किसी पिंड के \emph{बल आरेख}\index{बल आरेख} में केवल वही पिंड दिखता है,
और उस \emph{पर} लगता हर बल अपने \omterm{def:g10:inertia:force}{क्रिया बिंदु} से पैमाने
पर खींचा जाता है; पिंड \emph{द्वारा} लगाए गए बल उसमें कभी नहीं
आते।
\end{definition}

\begin{method}[किसी स्थैतिक स्थिति का विश्लेषण]\label{met:g10:inertia:method}
\begin{enumerate}
  \item जिस पिंड का अध्ययन करना है उसे चुनिए; \omterm{def:g10:relative-motion:frames}{भू-निर्देश तंत्र}
        में काम कीजिए।
  \item सूची बनाइए: हमेशा भार, और हर संपर्क के लिए एक
        \omterm{def:g10:inertia:contact}{संस्पर्श बल}।
  \item \omterm{def:g10:inertia:diagram}{बल आरेख} खींचिए।
  \item पिंड \omterm{def:g10:relative-motion:frame}{विराम} में है, इसलिए बल \omterm{def:g10:inertia:compensate}{संतुलनकारी} हैं
        (\cref{thm:g10:inertia:principle}): अज्ञात परिमाण आरेख से ही पढ़ लीजिए (एक विमा में बराबर
        लंबाइयाँ, दो विमाओं में घटक)।
\end{enumerate}
\end{method}

\begin{example}[मेज़ पर रखी पुस्तक]\label{ex:g10:inertia:book}
\qty{0.50}{kg} द्रव्यमान की एक पुस्तक मेज़ पर रखी है: भार $P = 0.50 \times 9.81 \approx \qty{4.9}{N}$ और
\omterm{def:g10:inertia:inventory}{अभिलंब प्रतिक्रिया} $\vect R$। \omterm{def:g10:relative-motion:frame}{विराम} में $\vect R$ $\vect P$ को
संतुलित करता है: \omterm{def:g10:universal-gravitation:weight}{ऊर्ध्वाधर}, ऊपर की ओर, $R = \qty{4.9}{N}$। एक और पुस्तक
रख दीजिए और $R$ उसी के हिसाब से बदल जाता है --- कोई आधार ठीक उतना ही
ज़ोर लगाता है जितना ज़रूरी हो।
\end{example}

\begin{omfigure}
\begin{tikzpicture}[scale=0.9, font=\small]
  \draw[fill=gray!15] (-2.6,2.2) rectangle (2.6,2.45);
  \draw (0,0) -- (-2.2,2.2);
  \draw (0,0) -- (2.2,2.2);
  \draw[dashed, gray] (0,0) -- (0,2.2);
  \draw[gray] (0,1.0) arc (90:135:1.0);
  \node[gray] at (-0.42,1.38) {$45^\circ$};
  \draw[fill=omExo!25] (0,-0.35) circle (0.35);
  \draw[-Stealth, omDef, very thick] (0,0) -- (-1.06,1.06)
    node[left] {$\vect T_1$};
  \draw[-Stealth, omDef, very thick] (0,0) -- (1.06,1.06)
    node[right] {$\vect T_2$};
  \draw[-Stealth, omThm, very thick] (0,0) -- (0,-1.5)
    node[right] {$\vect P$};
\end{tikzpicture}

{\small टँगा हुआ लैंप: दोनों तनावों के क्षैतिज हिस्से सममिति से
कट जाते हैं; उनके \omterm{def:g10:universal-gravitation:weight}{ऊर्ध्वाधर} हिस्से मिलकर भार उठाते हैं।}
\end{omfigure}

\begin{example}[दो तारों पर टँगा लैंप]\label{ex:g10:inertia:lamp}
\qty{3.0}{kg} द्रव्यमान ($P \approx \qty{29.4}{N}$) का एक लैंप दो तारों से टँगा है, और हर तार
\omterm{def:g10:universal-gravitation:weight}{ऊर्ध्वाधर} से $45^\circ$ पर है। सममिति से दोनों \omterm{def:g10:inertia:inventory}{तनाव} बराबर हैं
($T_1 = T_2 = T$) और उनके क्षैतिज घटक कट जाते हैं; \omterm{def:g10:universal-gravitation:weight}{ऊर्ध्वाधर} घटक, हर एक $T \cos 45^\circ$,
मिलकर भार उठाते हैं:
\[
2\,T \cos 45^\circ = P
\qquad\text{इसलिए}\qquad
T = \frac{29.4}{2 \times 0.707} \approx \qty{21}{N}
\]
--- यानी प्रति तार भार के आधे से भी अधिक: कोण बनाकर खींचने की
क़ीमत \omterm{def:g10:inertia:inventory}{तनाव} में चुकानी पड़ती है।
\end{example}

\section{अभ्यास}

\begin{exercise}[$\star$]\label{exo:g10:inertia:1}
इनमें से हर एक पर लगते बलों की सूची बनाइए और उन्हें संस्पर्श या दूरस्थ में
बाँटिए: (क) मेज़ पर \omterm{def:g10:relative-motion:frame}{विराम} में रखी पुस्तक; (ख) अपने तार से टँगा लैंप; (ग)
चिकनी बर्फ़ पर सरकती चकती; (घ) गिरता हुआ सेब (वायु का प्रतिरोध नगण्य)।
\end{exercise}

\begin{exercise}[$\star$]\label{exo:g10:inertia:2}
\qty{250}{g} की पुस्तक और \qty{60}{kg} के व्यक्ति का भार निकालिए। पैमाने
$\qty{1}{cm} \leftrightarrow \qty{1}{N}$ पर पुस्तक का तीर कितना लंबा होगा, और उस व्यक्ति के लिए यह पैमाना
बेकार क्यों है?
\end{exercise}

\begin{exercise}[$\star$]\label{exo:g10:inertia:3}
पैमाने $\qty{1}{cm} \leftrightarrow \qty{2}{N}$ वाले किसी आरेख में एक बल $A$ बिंदु पर लगा,
बाईं ओर इशारा करता \qty{3.5}{cm} लंबा क्षैतिज तीर है। उसकी तीनों पहचानें
बताइए।
\end{exercise}

\begin{exercise}[$\star$]\label{exo:g10:inertia:4}
\qty{1.2}{kg} द्रव्यमान का एक शब्दकोश किसी ताक़ पर रखा है। उसका \omterm{def:g10:inertia:diagram}{बल
आरेख} खींचिए और हर बल का परिमाण तथा दिशा बताइए।
\end{exercise}

\begin{exercise}[$\star$]\label{exo:g10:inertia:5}
सही या ग़लत --- ग़लत वालों को सुधारिए: (क) कोई पिंड आख़िर रुक ही जाता है,
जब तक कोई बल उसे धकेलता न रहे; (ख) \omterm{def:g10:inertia:uniform}{सरल रेखीय एकसमान
गति} के लिए किसी बिना कटे बल की ज़रूरत नहीं; (ग) \omterm{def:g10:relative-motion:frame}{विराम} में रखे
पिंड पर कोई बल नहीं लगता; (घ) यदि किसी पिंड का वेग बदलता है तो
उस पर लगते बल \omterm{def:g10:inertia:compensate}{संतुलनकारी} नहीं हैं।
\end{exercise}

\begin{exercise}[$\star\star$]\label{exo:g10:inertia:6}
छोड़े जाने के बाद एक चकती बर्फ़ पर \qty{24}{m} की दूरी \qty{3.0}{s} में, सीधे और
अचर चाल से तय करती है। उसकी चाल निकालिए; क्या उस पर लगते बल
\omterm{def:g10:inertia:compensate}{संतुलनकारी} हैं? कंक्रीट पर वही चकती मीटरों के भीतर रुक
जाती है: सूची में क्या बदला, और \cref{prop:g10:inertia:converse} क्या निष्कर्ष निकालता है?
\end{exercise}

\begin{exercise}[$\star\star$]\label{exo:g10:inertia:7}
\qty{70}{kg} का एक यात्री अचर \qty{1.5}{m/s} से ऊपर जाती लिफ़्ट में खड़ा है। यात्री
पर लगते बलों की सूची बनाइए; क्या वे \omterm{def:g10:inertia:compensate}{संतुलनकारी} हैं?
फ़र्श के बल का परिमाण बताइए। लिफ़्ट के \omterm{def:g10:relative-motion:frame}{विराम} में होने पर क्या
कुछ बदलेगा?
\end{exercise}

\begin{exercise}[$\star\star$]\label{exo:g10:inertia:8}
\qty{1.8}{kg} द्रव्यमान का एक लैंप एक ही \omterm{def:g10:universal-gravitation:weight}{ऊर्ध्वाधर} तार से \omterm{def:g10:relative-motion:frame}{विराम} में टँगा है।
तार का \omterm{def:g10:inertia:inventory}{तनाव} निकालिए। तार की जगह दो \emph{\omterm{def:g10:universal-gravitation:weight}{ऊर्ध्वाधर}} तार लगा दिए
जाते हैं जो बोझ बराबर बाँट लेते हैं: अब हर तार का \omterm{def:g10:inertia:inventory}{तनाव} कितना
है?
\end{exercise}

\begin{exercise}[$\star\star$]\label{exo:g10:inertia:9}
एक बस अचानक ब्रेक लगाती है और खड़ा यात्री आगे की ओर लुढ़क जाता है; और
अचर चाल पर मोड़ लेते समय वही यात्री बाहर की ओर झुक जाता है। किसी ``आगे
की ओर'' या ``बाहर की ओर'' बल को गढ़े बिना दोनों को समझाइए।
\end{exercise}

\begin{exercise}[$\star\star$]\label{exo:g10:inertia:10}
कुल \qty{85}{kg} द्रव्यमान का एक पैराशूट यात्री अचर \qty{5.0}{m/s} से
\omterm{def:g10:universal-gravitation:weight}{ऊर्ध्वाधर} नीचे उतरता है। \omterm{def:g10:inertia:diagram}{बल आरेख} खींचिए और छतरी पर
लगते वायु-प्रतिरोध की गणना कीजिए।
\end{exercise}

\begin{exercise}[$\star\star$]\label{exo:g10:inertia:11}
एक जल-स्कीबाज़ को अचर चाल पर सीधे खींचा जा रहा है; क्षैतिज रस्सी \qty{300}{N}
से खींचती है। स्कीयों पर पानी का क्षैतिज कर्षण कितना है (उचित ठहराइए)?
कौन-से दो \omterm{def:g10:universal-gravitation:weight}{ऊर्ध्वाधर} बल एक-दूसरे को संतुलित करते हैं?
\end{exercise}

\begin{exercise}[$\star\star\star$]\label{exo:g10:inertia:12}
\qty{2.4}{kg} द्रव्यमान का एक लैंप \omterm{def:g10:universal-gravitation:weight}{ऊर्ध्वाधर} से $40^\circ$ पर लगे दो तारों से टँगा
है ($\cos 40^\circ \approx 0.766$)। हर \omterm{def:g10:inertia:inventory}{तनाव} निकालिए। तारों के क्षैतिज होते जाने पर
\omterm{def:g10:inertia:inventory}{तनाव} क्यों बढ़ता है, और ठीक-ठीक क्षैतिज दो तार लैंप को कभी थाम
क्यों नहीं सकते?
\end{exercise}

\begin{exercise}[$\star\star\star$]\label{exo:g10:inertia:13}
कर्लिंग के एक पत्थर को धकेला जाता है, छोड़ा जाता है, वह बहुत धीरे-धीरे
मंद पड़ते हुए सरकता है, और फिर रुक जाता है। हर चरण (धक्का, सरकना, \omterm{def:g10:relative-motion:frame}{विराम})
के लिए \omterm{def:g10:inertia:diagram}{बल आरेख} खींचिए और बताइए कि बल
\omterm{def:g10:inertia:compensate}{संतुलनकारी} हैं या नहीं, और उसके साथ \cref{thm:g10:inertia:principle} या \cref{prop:g10:inertia:converse} का हवाला
दीजिए।
\end{exercise}

\begin{exercise}[$\star\star\star$]\label{exo:g10:inertia:14}
यान \emph{वॉयेजर 1} इंजन बंद किए हुए तारों के बीच के अंतरिक्ष में \qty{17}{km/s}
से लुढ़कता जा रहा है। अपनी चाल बनाए रखने के लिए उसे ईंधन क्यों नहीं
चाहिए, और अरस्तू क्या भविष्यवाणी करते? वह एक वर्ष में कितनी दूरी तय करता
है --- पहले \unit{km} में, फिर खगोलीय मात्रकों में ($\qty{1}{au} \approx
\qty{1.496e8}{km}$)? और
चाल बदले \emph{बिना} केवल दिशा बदलने के लिए भी उसे कोई प्रणोदक क्यों
चलाना पड़ता है?
\end{exercise}

\begin{exercise}[$\star\star\star$]\label{exo:g10:inertia:15}
तीन रस्सियाँ एक \omterm{def:g10:relative-motion:frame}{विराम} में रखे छल्ले को खींचती हैं। जिस जाली पर एक वर्ग
\qty{1}{N} का है, उस पर दो बल $\vect F_1\,(3, 2)$ और $\vect F_2\,(-1, 2)$ पढ़े जाते हैं।
तीसरे बल के घटक, परिमाण और दिशा निकालिए।
\end{exercise}

\section{समस्या: तीन अंकों में जड़त्व}

\begin{problem}[{सप्ताहांत समस्या --- हवाई अड्डे का चलपथ, बर्फ़ का मैदान
और पैराशूट: अचर वेग, संतुलित बलों की पहचान, और पैराशूट कर्षण नहीं बल्कि
चाल क्यों घटाता है}]
\label{pb:g10:inertia:1}
तीन दृश्य --- चलते पथ पर रखा एक सूटकेस, बर्फ़ के मैदान पर एक चकती, और
साँझ की हवा में उतरता एक पैराशूट यात्री --- और तीनों में बोलता एक ही
नियम।

\medskip
\noindent\textbf{भाग I --- हवाई अड्डे का चलपथ।}
\qty{23}{kg} का एक सूटकेस अचर चाल \qty{0.75}{m/s} से चलते पथ पर खड़ा है।
\begin{enumerate}
  \item भू-तंत्र में सूटकेस की गति क्या है? उस पर लगते बलों की
        सूची बनाइए।
  \item क्या ये बल \omterm{def:g10:inertia:compensate}{संतुलनकारी} हैं? उचित ठहराइए; दोनों
        परिमाण निकालिए।
  \item उसके \omterm{def:g10:inertia:diagram}{बल आरेख} की तुलना उसी सूटकेस के फ़र्श पर रखे होने
        वाले आरेख से कीजिए। यह तुलना कौन-सा गहरा कथन दिखाती है?
  \item चलपथ झटके से रुक जाता है। किसी आगे की ओर लगते बल को
        गढ़े बिना बताइए कि सूटकेस क्या करता है।
  \item उसकी मालकिन चलपथ के सापेक्ष \qty{1.2}{m/s} से चलती है। भू-तंत्र में
        उसकी चाल? क्या \emph{उस} पर लगते बल
        \omterm{def:g10:inertia:compensate}{संतुलनकारी} हैं?
\end{enumerate}

\medskip
\noindent\textbf{भाग II --- बर्फ़ का मैदान।}
\begin{enumerate}[resume]
  \item \qty{160}{g} की एक चकती अचर \qty{8.0}{m/s} से सीधी सरकती है। परिमाणों सहित
        उसका \omterm{def:g10:inertia:diagram}{बल आरेख} खींचिए।
  \item अरस्तू का दावा है कि चकती को कोई चलाने वाला चाहिए। मैदान क्या
        जवाब देता है, और रोज़मर्रा की सतहें उनका पक्ष क्यों लेती दिखती
        हैं?
  \item चकती किसी खुरदरे टुकड़े से गुज़रती है और धीमी पड़ जाती है। वहाँ
        के बलों के बारे में क्या निष्कर्ष निकलेगा? दोषी
        बल कौन-सा है?
  \item एक स्केटर अचर चाल पर सहज मोड़ लेती है। क्या उस पर लगते
        बल \omterm{def:g10:inertia:compensate}{संतुलनकारी} हैं? सावधान: वेग किस तरह की राशि
        है?
  \item आरेख को गति से मिलाइए: (क) \omterm{def:g10:inertia:compensate}{संतुलनकारी बल}; (ख) $\vect v$
        के उलटा एक अतिरिक्त बल; (ग) $\vect v$ के लंबवत एक
        अतिरिक्त बल --- (1) सीधे चलते हुए धीमा पड़ना; (2) अचर
        चाल पर मुड़ना; (3) \omterm{def:g10:inertia:uniform}{सरल रेखीय एकसमान गति}।
\end{enumerate}

\medskip
\noindent\textbf{भाग III --- पैराशूट।}
सामान समेत छलाँग लगाने वाले का कुल द्रव्यमान \qty{80}{kg} है।
\begin{enumerate}[resume]
  \item विमान छोड़ते ही \omterm{def:g10:universal-gravitation:weight}{ऊर्ध्वाधर} चाल और कर्षण, दोनों अभी शून्य हैं।
        बल निकालिए; क्या वे \omterm{def:g10:inertia:compensate}{संतुलनकारी} हैं? आगे क्या
        होना चाहिए?
  \item कर्षण चाल के साथ बढ़ता है। जब तक कर्षण भार से $<$ है
        तब तक की गिरावट बताइए।
  \item छतरी खुलने से पहले छलाँग लगाने वाला स्थिर \qty{50}{m/s} (``सीमांत
        चाल'') तक पहुँच जाता है। कर्षण निकालिए। गति किस तरह की है, और वह
        ``भेस बदले जड़त्व'' क्यों है?
  \item छतरी खुलती है: अब कर्षण भार से कहीं अधिक है। वेग का
        क्या होता है? (छलाँग लगाने वाला पूरे समय \emph{नीचे} ही चलता
        है।)
  \item एक नई स्थिर \qty{5.0}{m/s} बन जाती है। कर्षण निकालिए; प्रश्न 13 से
        तुलना कीजिए।
  \item और अंतिम चोट: पैराशूट ने क्या बदला --- सीमांत कर्षण, या वह चाल
        जिस पर कर्षण भार तक पहुँच जाता है?
\end{enumerate}

\medskip
\noindent\textbf{भाग IV --- फ़ैसले।}
\begin{enumerate}[resume]
  \item पूरा कीजिए और उचित ठहराइए: ``\omterm{def:g10:relative-motion:frame}{विराम} में या \omterm{def:g10:inertia:uniform}{सरल रेखीय
        एकसमान गति} में $\iff$ \dots''। \omterm{def:g10:relative-motion:frame}{विराम} को अलग नियम की ज़रूरत
        क्यों नहीं?
  \item किस अंक में अरस्तू का ``गति को चलाने वाला चाहिए'' सबसे खुलकर
        धराशायी होता, और गैलीलियो किस ओर इशारा करते?
  \item किसी परीक्षा की उत्तर-पुस्तिका में लिखा है: ``पैराशूट यात्री अचर
        चाल से गिरता है, इसलिए उस पर कोई बल नहीं लगता''। इसे एक
        पंक्ति में सुधारिए।
  \item समापन: हर अंक के लिए उन बलों का युग्म बताइए जो
        संतुलित होता है, और वह इकलौता नियम बताइए जो इस समस्या
        में बीस बार काम आया।
\end{enumerate}
\end{problem}
